\section{Scalar bounds for the Gaussian construction}
\label{sec:arithmetic}

The elementary inequalities and scalar error budgets below are the numerical
inputs to the interpolation and sign arguments in
Sections~\ref{sec:interpolation} and~\ref{sec:signs}.  We give rational
proofs of these bounds.  The supplied program
\path{verification/arithmetic_bounds.py} checks these scalar budgets and the
numerical error-propagation majorants in Appendix~\ref{app:floor-certificate}, using $512$-bit Arb balls, exact rational
arithmetic, and exact integer samples of the rounded exponential routine.
It does not implement the degree-$2000$ rational exponential enclosures described below.

The finite matrix certificate has a separate verification.  The program
\path{verification/numeric_balls.py} encloses the exact quadrature arrays at
$256$-bit precision and tests every strict comparison in
Tables~\ref{table:power-bounds}--\ref{table:sign-bounds}.  This is the finite
route used in Proposition~\ref{prop:finite-certificate}; its relation to the
exact spectral integrals is supplied by Lemma~\ref{lem:quadrature-error}.

Appendix~\ref{app:floor-certificate} gives a separate conditional
floor-rounded rational procedure, with all its error estimates and finite
loops.  That full matrix construction is not executed by the supplied
programs.  Commands and dependencies for the two programs are listed in
\path{verification/README.md}.

\subsection{Rational enclosures of elementary constants}

Every terminating decimal specifying a parameter or comparison threshold is
an exact rational number.  Set $\epsilon=10^{-80}$ and define
\begingroup\small
\begin{equation}\label{eq:rational-pi-b}
\begin{split}
 p_*={}&3.14159265358979323846264338327950288419716939937510582097494459230781640628620899,\\
 b_*={}&0.86602540378443864676372317075293618347140262690519031402790348972596650845440001.
\end{split}
\end{equation}
\endgroup

\begin{lemma}[Decimal enclosures]\label{lem:rational-pi-b}
For $b=\sqrt3/2$, these rationals satisfy
\begin{equation}\label{eq:rational-pi-b-error}
 p_*<\pi<p_*+\epsilon,\qquad b_*<b<b_*+\epsilon.
\end{equation}
\end{lemma}
\begin{proof}
Put
\[
 A(x)=\sum_{j=0}^{65}\frac{(-1)^jx^{2j+1}}{2j+1}.
\]
For $0<x<1$, integrate the finite geometric expansion of $1/(1+x^2)$.
Its remainder has the sign and magnitude giving
$A(x)<\arctan x<A(x)+x^{133}/133$.  The tangent addition formula gives
$\tan(4\arctan(1/5))=120/119$; subtracting an angle with tangent $1/239$
gives tangent one, and the resulting angle is in $(0,\pi/2)$.  Therefore
$\pi/4=4\arctan(1/5)-\arctan(1/239)$.  The two required rational comparisons
are
\[
\begin{split}
 p_*&<16A(1/5)-4\bigl(A(1/239)+239^{-133}/133\bigr),\\
 16\bigl(A(1/5)+5^{-133}/133\bigr)-4A(1/239)&<p_*+\epsilon.
\end{split}
\]
They follow by clearing positive denominators.  Likewise, the second bracket
is exactly the pair of integer-square comparisons
$b_*^2<3/4<(b_*+\epsilon)^2$.
\end{proof}

The threshold in Lemma~\ref{lem:gaussian-duality} needs only a shorter
rational comparison.  Squaring gives $\sqrt3<1733/1000$.  Using the
Machin identity proved above, the alternating arctangent series gives
\[
 \pi>16\left(\frac15-\frac1{3\cdot5^3}\right)-\frac4{239}
 =\frac{281476}{89625}>\frac{157}{50}.
\]
Consequently
\[
 \pi(2/\sqrt3-2/5)
 >\frac{157}{50}\left(\frac{2000}{1733}-\frac25\right)
 =\frac{1025838}{433250}>\frac{59}{25}=2.36.
\]

We will also use a rational enclosure of a real exponential.  For rational
$a\ge0$ and an integer $q\ge0$, put $P_q(a)=\sum_{j=0}^qa^j/j!$.
If $a>0$ and $q+2>a$, comparison of successive terms in the positive tail
gives
\begin{equation}\label{eq:scalar-rational-enclosure}
\begin{gathered}
 P_q(a)<e^a<
 P_q(a)+\frac{a^{q+1}}{(q+1)!}\frac1{1-a/(q+2)},\\
 \frac1{P_q(a)+\frac{a^{q+1}}{(q+1)!}\frac1{1-a/(q+2)}}
 <e^{-a}<\frac1{P_q(a)}.
\end{gathered}
\end{equation}
Indeed, the first omitted term is $a^{q+1}/(q+1)!$, and every subsequent
ratio is at most $a/(q+2)<1$.  The lower bound is the positive partial sum,
and reciprocation reverses the two positive bounds.  At $a=0$ both
exponentials equal one.  Lemma~\ref{lem:quadrature-error} uses only the
$a=1,q=6$ instance to prove $e<2719/1000$; its three final numerical
allowances have already been established there.

\subsection{Scalar bounds for interpolation and signs}

We now evaluate the scalar majorants derived in
Sections~\ref{sec:interpolation} and~\ref{sec:signs}.  The first is the tail
bound from Lemma~\ref{lem:kernel-bounds}, for a positive integer $r$:
\[
 U_r(w,c)=
 \frac{800(w+c)e^{-.70r}}{1-e^{-.70}}
 +\frac{5.3(2.54w/r+2c)e^{-.187r}}{1-e^{-.187}}.
\]
In each application, $w$ bounds the sum of the absolute node coefficients
and $c$ bounds the absolute constant coordinate, summed over the input lists
when there are two.  Lemma~\ref{lem:kernel-bounds} then bounds the sum of
the absolute transformed jet coordinates over all output nodes $n\ge r$
by $U_r(w,c)$.

\Needspace{8\baselineskip}
The following strict outward bounds collect the numerical inputs for
interpolation and the first sign estimates:
\begin{equation}\label{eq:scalar-error-budgets}
\begin{array}{c|c|c}
\text{quantity}&\text{outward upper bound}&\text{allowance used}\\ \hline
 U_{168}(1,0)&1.0670\cdot10^{-14}&1.1\cdot10^{-14}\\
 U_{168}(5,.02)&8.1575\cdot10^{-14}&8.3\cdot10^{-14}-10^{-170}\\
 U_{91}(5,.02)&2.2711\cdot10^{-7}&2.4\cdot10^{-7}\\
 \displaystyle\frac{1+(1+2.12)/2.36}{1.75}
       \frac{e^{-2.36\cdot11}}{1-e^{-2.36}}
       &7.7915\cdot10^{-12}&8\cdot10^{-12}\\[2pt]
 \displaystyle(2.3+5.6/2.36)
       \sum_{n\in\{1,3,4,7,9\}}e^{-2.36(n-1)}
       &4.7185&4.8\\[2pt]
 \displaystyle\frac{36}{2}(2\cdot10^{-8})3(3\pi)^2
       &9.5933\cdot10^{-5}&.00011
\end{array}
\end{equation}
The first tail residual has a deliberately finer margin:
$8.3\cdot10^{-14}-8.1575\cdot10^{-14}
=1.425\cdot10^{-15}>10^{-170}$.  For the high-node list, the finer
$2.2711\cdot10^{-7}$ bound is retained before the target jets are added.

\paragraph{Local target tails and interpolation arithmetic.}
Write
\[
 c_{\mathrm{jet}}=\frac{1+(1+2.12)/2.36}{1.75}
 =\frac{548}{413},\qquad \rho=e^{-2.36}.
\]
The target-jet estimate \eqref{eq:target-jet-bound}, monotonicity in
$k\ge2.36$, and overcounting by all integers show that the target cost
from nodes $n\ge r+1$ is at most
$c_{\mathrm{jet}}\rho^r/(1-\rho)$.  A small rational check suffices even
for the very distant tail:
\[
 P_7(59/25)>1+2.36+2.78+2.19+1.29+.61+.239+.08
 =10.549>21/2.
\]
Each termwise lower comparison is rational.  Thus
$\rho<2/21$.  Also
\[
 \frac{548}{413}\frac{21}{19}<\frac32,\qquad
 (21/20)^{168}>(21/20)^{160}>
 (3/2)^{16}>150.
\]
The middle inequality follows from
$(1+1/20)^{10}>1+10/20=3/2$, and the last follows by comparing
$3^{16}$ with $150\cdot2^{16}$.  Therefore the target tail above $168$
satisfies the explicit local bound
\begin{equation}\label{eq:target-tail-168}
 c_{\mathrm{jet}}\frac{\rho^{168}}{1-\rho}
 <\frac{548}{413}\frac{21}{19}(2/21)^{168}
 <\frac{3}{2}\frac{10^{-168}}{150}=10^{-170}.
\end{equation}
Combining this with the finer $U_{168}(5,.02)$ row of
\eqref{eq:scalar-error-budgets} proves the tail allowance
$U_{168}(5,.02)+10^{-170}<8.3\cdot10^{-14}$.
This direct rational bound handles the exponent
$2.36\cdot168=396.48$, outside the range of exponents used in the
degree-$2000$ comparisons below.

The same local estimate for the first high node $91$ gives
\begin{equation}\label{eq:target-tail-91}
 c_{\mathrm{jet}}\frac{e^{-2.36\cdot90}}{1-e^{-2.36}}
 <\frac32\,10^{-90}<10^{-89}.
\end{equation}
Thus the transformed and target contributions to the high-list norm obey
$U_{91}(5,.02)+10^{-89}<2.4\cdot10^{-7}$.  The finer $U_{91}$ bound
leaves the margin needed to include the target jets.

The same $\rho<2/21$ gives a direct bound for the first five target jets:
\[
 \frac{548}{413}\left(1+(2/21)^2+(2/21)^3+(2/21)^6+(2/21)^8\right)
 <1.35.
\]
The reference-list bound in \eqref{eq:scalar-error-budgets} includes the
extras, so their combined norm is below $4.8$.  Their constant coordinates
also retain a stronger bound.  For $k\ge2.36$, the first expression below
is negative and the second positive, whence
\[
 |-.0074+.0115/k|+|.0096-.017/k|
 =.017-.0285/k<.02.
\]
The first-five jet bound and the quadrature-operator bound imply
$1.35+(4000+2\cdot10^{-27})4.8<20000$ for the combined data norm
used in the geometric residual.

For clarity, the remaining local inverse and residual arithmetic is
\[
\begin{gathered}
 \frac{36}{1-.004^8}<37,\qquad
 \frac{37}{1-37(2\cdot10^{-27})}<40,\\
 (1.1\cdot10^{-14})(1+40\cdot4000)
 =1.760011\cdot10^{-9}<1,\\
 20000(.004)^8=1.31072\cdot10^{-15}<1.4\cdot10^{-15}.
\end{gathered}
\]
The finite residual is consequently bounded by
$8\cdot10^{-12}+10^{-24}+1.4\cdot10^{-15}<10^{-11}$ after its three
terms have been derived in Section~\ref{sec:interpolation}.
For the correction inequalities \eqref{eq:correction-system}, put
$\beta=1.1\cdot10^{-14}$.  Their rational solution is bounded by
\[
 u_{\mathcal T}\le
 \frac{8.3\cdot10^{-14}+40\beta\cdot10^{-11}}
 {1-\beta(1+40\cdot4000)},\qquad
 u_{\mathcal F}\le40(10^{-11}+4000u_{\mathcal T}).
\]
The denominator is positive by the preceding Schur bound, and the two
right sides give
$u_{\mathcal F}+u_{\mathcal T}<1.3681\cdot10^{-8}<2\cdot10^{-8}$.
Together with the operator-derived inequalities
\eqref{eq:correction-system}, this proves the correction bound.  The
reference norm then gives a combined
exact-list norm below $4.8+2\cdot10^{-8}<5$, with the fixed constant
coordinates still bounded by $.02$.

\paragraph{Half-gap and tail sign arithmetic.}
For the Taylor remainder in \eqref{eq:half-gap-series-tail}, put
\[
 \mathcal R(m,Y,\nu)=
 \frac{Y^{-\nu}36\cdot5}
 {(41+\nu)!\bigl(1-3\pi Y/(42+\nu)\bigr)}
 \left((\pi Y)^{41+\nu}
 +3\sum_{j=1}^6e^{-\pi\Ima z(t_j)m}
       (\pi|z(t_{j-1})|Y)^{41+\nu}\right).
\]
\begin{samepage}
The denominators are positive for the four tuples below; for example,
$3\pi Y/(42+\nu)\le(9\pi/2)/42<1$.  The sum of six nonnegative
candidates dominates the maximum in the analytic remainder.  The strict
outward bounds are
\[
\begin{array}{c|c}
(m,Y,\nu)&\mathcal R(m,Y,\nu)\text{ is below}\\ \hline
(0,1/2,0)&1.8\cdot10^{-19}\\
(1,1,2)&1.0\cdot10^{-11}\\
(3,1,2)&3.0\cdot10^{-16}\\
(4,3/2,2)&2.0\cdot10^{-11}.
\end{array}
\]
\end{samepage}
The included-moment quadrature cost is separately at most
$41\cdot4\cdot5\cdot10^{-31}=8.2\cdot10^{-29}$.
Adding it to each displayed remainder still gives less than $10^{-8}$.
As proved after \eqref{eq:half-gap-series-tail}, the bound decreases when
$m$ increases or $Y$ decreases, so these four tuples cover every half-gap.

The two target-support Taylor tails satisfy
\[
 \frac{e^6}{6}\frac{3^{41}}{41!}\frac1{1-3/42}
 <8\cdot10^{-29}<10^{-6},\qquad
 6\frac{6^{41}}{43!}\frac1{1-6/44}
 <10^{-20}<10^{-6}.
\]
These estimates are applied after the parameter monotonicity and signed
half-gap cases in Lemma~\ref{lem:exponential-support}; in particular,
$K\le6$ also in the final unbounded parameter box.

At $s_0=89.5$, the four terms in the tail second-derivative majorant obey
\begin{align*}
 &36\pi^2(2.4\cdot10^{-7})
 +3(3\pi)^2\,36\cdot5\,e^{-.70s_0}\\
 &\quad
 +1.25(1.32\pi)^2(2.54\cdot5/s_0+.04)e^{-.187s_0}
 +2.36e^{-2.36(s_0-1)}
 <8.5485\cdot10^{-5}<.0002.
\end{align*}
The last term is the supremum of $ke^{-k(s_0-1)}$ for $k\ge2.36$:
its logarithmic derivative is $1/k-(s_0-1)<0$ there.
The smallest endpoint ratio for the divided periodic product satisfies
$12-8\sqrt2>.6862>.68$.  The midpoint identities and logarithmic concavity in
Lemma~\ref{lem:sine-product-barrier} extend this endpoint bound to every
nearest-node half-gap.

The final numerical implications are elementary once their separate error
terms and barriers have been proved:
\[
 .001-.00011-.000001-.000001=.000888>0,\qquad
 .0021-2\cdot10^{-8}>.002,\qquad
 .68\cdot.002>.0001.
\]
These give, respectively, the bounded-range quotient margin, the exact
low-node rational margin, and the comparison with the tail remainder.

\paragraph{A nonmatrix rational route for these scalars.}
For the three $U_r$ evaluations and the reference-list, omitted-jet,
deleted-quotient, four half-gap-remainder, two target-support-tail,
second-derivative, and product
endpoint comparisons above, use \eqref{eq:scalar-rational-enclosure} with
$q=2000$.  Square roots, when present, are enclosed between positive
rationals of width $10^{-80}$ whose squares straddle the radicand.
Explicitly, write a radicand as $A_{\mathrm{num}}/A_{\mathrm{den}}$ with
positive integers $A_{\mathrm{num}},A_{\mathrm{den}}$.  Integer search over
$0\le m<3\cdot10^{80}$ finds $m$ with
$m^2A_{\mathrm{den}}\le10^{160}A_{\mathrm{num}}<(m+1)^2A_{\mathrm{den}}$.
Then $m/10^{80}\le\sqrt{A_{\mathrm{num}}/A_{\mathrm{den}}}<(m+1)/10^{80}$;
equality at the lower endpoint is kept exact.  The radicands used here exceed
one, so both bracket endpoints are positive.
The tail and interpolation comparisons need only rational arithmetic and
exponential enclosures, with positive geometric denominators.
The deleted-quotient, half-gap-remainder, and second-derivative comparisons
also use the $\pi$ output of Lemma~\ref{lem:rational-pi-b}; they do not
need its $b$ output.  For the half-gap remainders,
\[
 \Ima z(t)=\frac{Bh}{t^2+h^2}-h,\qquad
 |z(t)|^2=h^2+\frac{B(B-2h^2)}{t^2+h^2}
\]
at rational $t=t_j$, so only the latter positive rational radicand needs
a square-root bracket.  For $0\le t\le1$,
$1<4/25+304/261\le |z(t)|^2<9$, by this formula and $|z|<3$.
The endpoint ratio uses a bracket for $\sqrt2$, whose radicand also lies
between one and nine.
Thus each of these scalar comparisons requires no matrix construction.

All exponential magnitudes in this list are below $250$: the largest
explicit target phase is $2.36(s_0-1)=208.86$, the largest tail phase
is $.70\cdot168=117.6$, and the four half-gap phases are below
$\pi\cdot3\cdot4<38$ by $|z|<3$.  Finite rational evaluation with
$q=2000$ and the indicated root brackets gives the printed outward
margins.  The two more distant target-jet tails instead use the local
$P_7(59/25)$ calculation above; the quadrature proof uses $q=6$.

\section{Periodic approximation for the field energy}
\label{app:normalized-periodic-approximation}

Fix $s$ and the exponent $p$ as in Section~\ref{sec:renormalized}.
Throughout this appendix we suppress the subscript $s$ in
$g_s,w_s,\kappa_s,f_{s,\eta},\mathcal W_{s,\eta},\mathcal W_s$, and $V_s$,
retaining the same order of limits.
Thus $L=-\operatorname{div}(w\nabla\cdot)$,
$Lg=\kappa\delta_0$, and $\kappa=2\pi$ for the planar logarithm,
whereas $\kappa=4\pi$ for $g(X)=|X|^{-s}$, $0<s<2$, with
$k=1$, $\gamma=s-1$, and $w=|y|^\gamma$ on the full extension
space; for the logarithm set $k=0$, $\gamma=0$, and $w=1$.
A compatible field also
satisfies $wE\in L^1_{\rm loc}$; its equation is understood through
the finite weak integrals $\int wE\cdot\nabla\varphi$.
This requirement is automatic whenever a truncation has finite local
weighted energy.  Write $K_R=[-R/2,R/2]^2$ and
\[
 Q_\eta(E,R)=\int_{K_R\times\mathbb R^k}w|E_\eta|^2,
 \qquad
 F_\eta=\int_{\mathbb R^2}f_\eta(x,0)\,dx
 =\begin{cases}\pi\eta^2/2,&s=0,\\
 \pi s\eta^{2-s}/(2-s),&s>0.
 \end{cases}
\]

\begin{proposition}[Normalized cubic approximation]
\label{prop:normalized-cubic-approximation}
Let $0\le s<2$, and fix
$1<p<2$ for $s=0$, or
$1<p<\min(2,2/s,3/(s+1))$ for $s>0$.
Let $\mathcal A_1$ be the compatible gradient fields in this fixed
$L^p_{\rm loc}$ class, with simple locally finite configurations and
unit background.  Then the limit
\[
 \mathcal W(E)=\lim_{\eta\downarrow0}
 \left[\limsup_{R\to\infty}R^{-2}Q_\eta(E,R)
                    -\kappa g(\eta)\right]
\]
exists in $\mathbb R\cup\{+\infty\}$ and is bounded below uniformly
over $\mathcal A_1$.  Its infimum $m$ is finite.
There are integers $R_j\to\infty$ and simple configurations
$\mathcal C_j$, periodic under $2R_j\mathbb Z^2$, with
$4R_j^2$ points per period cell, whose canonical periodic values satisfy
\[
                   V(\mathcal C_j)\longrightarrow m.
\]
The approximating fields may therefore be chosen to be canonical
periodic gradients in the same fixed-$p$ class.
\end{proposition}

The proof has three inputs: cutoff control, a stationary minimizing law,
and a screening construction. The following lemmas establish them before
we reflect and project the screened fields to prove the proposition.

A cutoff-energy estimate alone does not control the renormalized energy:
removing the cutoff can expose arbitrarily close pairs.  We first obtain a
stationary law supported on minimizing fields whose expected close-pair
defect tends to zero with the cutoff.  We then screen one realization
without creating new close pairs, and reflect it to obtain a periodic
configuration.  The periodic defect identity below converts its cutoff
bound into a bound for the renormalized energy.

The screening fields need not be gradients.  At that stage we preserve
the divergence equation, obtain zero exterior normal flux, and bound the
cutoff energy.  After reflection, projection onto the canonical periodic
gradient preserves the charges and decreases that energy.
For cutoff estimates and compactness we temporarily allow integer
multiplicities; the stationary-law lemma will rule out collisions before
using the infimum over simple configurations.
Let $B_\eta$ be the ball of radius $\eta$ centered at zero in
$\R^{2+k}$. The positive unit sphere measure is
\[
 d\delta_0^{(\eta)}
   =\kappa^{-1}[-g'(\eta)]w\,d\sigma_{\partial B_\eta},
 \qquad Lf_\eta=\kappa(\delta_0-\delta_0^{(\eta)}).
\]
The measure $\delta_p^{(\eta)}$ is its translate to $(p,0)$.
The sign is fixed by the decreasing kernel. 
Write $\nu=\sum_pN_p\delta_p$, $\nu_\eta=\sum_pN_p\delta_p^{(\eta)}$,
and $N_B=\nu(B)$.
Fix the geometric constants
\[
 a_0=\frac14\sqrt{\frac23},\qquad a_1=4\sqrt2,
 \qquad \eta_0=\min\left\{\frac18,\frac{a_0}{16}\right\}.
\]
These constants are independent of the cutoffs, the realization,
$\varepsilon$, and $R$.  Indeed, the planar subdivision in
Petrache and Serfaty~\citep[Lemma~6.3]{PetracheSerfaty2017} gives sidelengths in
$[\frac14\widetilde m^{-1/2},4\widetilde m^{-1/2}]$;
for $\widetilde m\in[1/2,3/2]$ this interval is contained in
$[a_0,a_1]$.  Fix once and for all a dyadic auxiliary cutoff
$\eta'\in(0,\eta_0)$ for the local count estimates, and include it
among the cutoffs controlled by the stationary law.  The cutoffs used for
stationary approximation will run through dyadic values $0<\eta<\eta'$.


\subsection{Cutoff control and spatial density}

\begin{lemma}[Cutoff control]\label{lem:field-cutoff-control}
For compatible gradients in the prescribed fixed-$p$ class, allowing
integer multiplicities, the cutoff limit $\mathcal W(E)$ exists in
$\R\cup\{+\infty\}$. There is a constant $C_s$ independent of the field
such that
\[
 -C_s\le\mathcal W_\eta(E)\le\mathcal W(E)+2\kappa F_\eta
 \qquad(0<\eta<\eta_0).
\]
If $Q_\eta(E,R)=O(R^2)$ at one fixed cutoff, $Q_\rho(E,R)=O(R^2)$ holds at every
fixed cutoff $\rho$ and $\nu(K_R)/R^2\to1$. In particular, the infimum $m$ in
Proposition~\ref{prop:normalized-cubic-approximation} is a finite real number.
\end{lemma}

\begin{proof}
Testing the
smeared field equation against a nonnegative smooth function equal to
one near all spheres centered in a fixed unit cell gives
\begin{equation}
 N_B^2\le C_{B,\eta}\left(1+
          \int_{B'\times[-h,h]^k}w|E_\eta|^2\right).
 \label{eq:normalized-local-count}
\end{equation}
Here $B'$ and $h$ are fixed enlargements.  These enlargements have
bounded overlap over a unit-cell tiling.  Also
\[
 \int w|\nabla(f_\alpha-f_\eta)|^2
       =\kappa[g(\alpha)-g(\eta)],\qquad 0<\alpha<\eta.
\]
Consequently, if $Q_\alpha(E,R)=O(R^2)$ at one fixed cutoff, then
$Q_\eta(E,R)=O(R^2)$ at every fixed cutoff: apply
\eqref{eq:normalized-local-count} to the finite-range difference
$\sum_pN_p\nabla(f_\alpha-f_\eta)(\cdot-p)$.

The smeared Gauss equation also gives
$\nu(K_R)/R^2\to1$ whenever one cutoff has volume-order energy.
For completeness, choose a lateral boundary at distance at most one
from $\partial K_R$ with squared flux energy $O(R^2)$.
For $k=1$, choose height $L\in[1,\sqrt R]$ with top and bottom
squared energy $O(R^{3/2})$.  Weighted Cauchy--Schwarz bounds the
total flux by
\[
 C L^{(\gamma+1)/2}R^{3/2}
       +C L^{\gamma/2}R^{7/4}=o(R^2),
\]
since $-1<\gamma<1$ (for negative $\gamma$, use $L^{\gamma/2}\le1$).
The smeared count lies between the counts in squares whose sidelengths
differ from $R$ by a fixed constant.  Applying these inequalities at
the correspondingly shifted radii proves full-radius density
convergence.  In the logarithmic case only the lateral term is needed.
This is the argument of Petrache and Serfaty~\citep[Lemma~2.1]{PetracheSerfaty2017}, with
the actual flux in the weak equation.

Put $S=\sum_pN_p(f_\alpha-f_\eta)(\cdot-p)\le0$.
For a horizontal cutoff $\chi$, integration by parts gives
\begin{align*}
 \int\chi w(|E_\eta|^2-|E_\alpha|^2)
  &=\kappa\int\chi S(\nu_\eta+\nu_\alpha-2\delta_{\mathbb R^2})\\
  &\quad-\int S w(E_\alpha+E_\eta)\cdot\nabla\chi.
\end{align*}
Distinct-charge terms are nonpositive.  An interior self term is
$-\kappa[g(\alpha)-g(\eta)]N_p^2$, bounded above by the same
expression with $N_p$ in place of $N_p^2$.  The background contributes
at most $2\kappa F_\eta$ per particle.  The commutator and transition
terms are bounded by the sum of the two energies and the counts in a
fixed-width shell, using \eqref{eq:normalized-local-count}.

To remove that shell, choose $r_j$ realizing the $\eta$-limsup.
Averaging both fixed-cutoff shell energies over
$t\in[r_j+b_0,r_j+\sqrt{r_j}]$, where $b_0$ exceeds twice the
shell width, gives a radius $t_j$ at which their sum is
$O(r_j^{3/2})=o(r_j^2)$.
Since $t_j/r_j\to1$ and $Q_\eta$ is nonnegative and increasing,
\[
 \frac{Q_\eta(E,t_j)}{t_j^2}
   \ge\frac{r_j^2}{t_j^2}\frac{Q_\eta(E,r_j)}{r_j^2}.
\]
Thus these radii still realize the $\eta$-limsup.  Only an upper
limsup bound is needed for $Q_\alpha(E,t_j)/t_j^2$.
The count shell is $o(t_j^2)$ by the density conclusion.  Hence
\begin{equation}
 \mathcal W_\eta(E)\le\mathcal W_\alpha(E)+2\kappa F_\eta,
 \qquad 0<\alpha<\eta<\eta_0.
 \label{eq:normalized-almost-monotonicity}
\end{equation}
These selected radii suffice for the upper-limit comparison.
Comparing with one fixed cutoff proves a uniform lower bound;
then \eqref{eq:normalized-almost-monotonicity} proves existence of the
full cutoff limit and
\begin{equation}
 -C_s\le\mathcal W_\eta(E)\le\mathcal W(E)+2\kappa F_\eta.
 \label{eq:normalized-cutoff-bound}
\end{equation}
If every cutoff has infinite energy, the conclusion is immediate.
These arguments also hold with integer multiplicities.
A canonical unit-lattice field has finite renormalized value, so
$-\infty<m<\infty$.

\end{proof}

\subsection{A minimizing stationary law and its close-pair defect}

The cutoff estimate does not exclude rare clusters. We now use translation
averages to obtain a stationary law for which their contribution can be
measured exactly. Fix the unit square $Q=[-1/2,1/2)^2$ and a decreasing
dyadic sequence $\eta_\ell\downarrow0$ beginning with $\eta_1=\eta'$.
For distinct locations define the nonnegative kernel
\[
 q_\eta(p-q)=f_\eta(p-q,0)
             +\int f_\eta(X-(p,0))\,d\delta_q^{(\eta)}(X).
\]
It vanishes when $|p-q|\ge2\eta$.
Anchor each ordered pair at its first point and define the measure
\[
 \mu_\eta=\kappa\sum_{p\in\mathcal C}
          \left(\sum_{q\in\mathcal C\setminus\{p\}}q_\eta(p-q)\right)\delta_p.
\]

\begin{lemma}[Stationary minimizers with vanishing defect]
\label{lem:stationary-field-minimizers}
There is a horizontally stationary probability law $P$ on compatible
simple gradient fields in the fixed-$p$ class, supported on
$\mathcal W(E)=m$. For the selected dyadic cutoffs, put
\[
 a_\eta=\mathbb E_P\mathcal W_\eta(E),\qquad
 D_\eta=\mathbb E_P\mu_\eta(Q).
\]
Then $a_\eta\to m$ and $0\le D_\eta\le m-a_\eta+2\kappa F_\eta\to0$.
For each selected $\eta$, one can choose a realization for which the
spatial averages of every selected cutoff energy have finite limits, the density
$d_\eta=\lim_{R\to\infty}R^{-2}\mu_\eta(K_R)$ satisfies
$d_\eta\le2D_\eta$, and, when $k=1$, for every fixed $\varepsilon>0$,
\[
 \frac1{\varepsilon^4R^2}
 \int_{K_R\times\{|y|>\varepsilon^2R/4\}}w|E|^2\longrightarrow0.
\]
\end{lemma}

\begin{proof}
Choose simple fields $E^j$ with $\mathcal W(E^j)\le m+1/j$.
For each $j$ choose one radius $S_j$, with $S_j/j\to\infty$,
large enough that the upper limsup estimates for $E^j$ hold
simultaneously at $\eta_1,\ldots,\eta_j$ on $K_{S_j+2j}$.
This enlargement handles all observation windows contained in the
centered exhaustion square $K_{2j}$.
This is possible because there are only finitely many estimates at
each stage.  Let $P_j$ be the law of
$(E^j(\cdot+(x,0)),\nu_j(\cdot+x))$ when $x$ is uniform in $K_{S_j}$.
For the unit square $Q$ fixed above, Tonelli's theorem and the
enlargement by $2j$ give the limiting bounds
\begin{equation}
 \limsup_j\int\Gamma_{\eta_\ell}\,dP_j
 \le m+2\kappa F_{\eta_\ell}+\kappa g(\eta_\ell),
 \qquad
 \Gamma_\eta(E)=\int_{Q\times\mathbb R^k}w|E_\eta|^2.
 \label{eq:normalized-orbit-bound}
\end{equation}
The analogous bounds hold on every fixed base square.

Here are the compactness details.  Use $H^{-r}_{\rm loc}$ for $E$,
with $r>(2+k)/2+1$, together with the vague topology for
the positive charge measure.  Bounded local $L^p$ sets are relatively
compact in the first space.  Equation
\eqref{eq:normalized-local-count} bounds second moments of local
counts.  Weighted H\"older bounds $E_\eta$ in local unweighted $L^p$
when $p<\min(2,2/s)$; the singular correction belongs to local $L^p$
when $p<3/(s+1)$.  For the logarithm both requirements reduce to
$p<2$.  Thus Markov's inequality and a countable exhaustion give
tightness.  Allow multiplicities in the resulting integer point
measure.

On sets with bounded local counts, vague convergence identifies every
fixed singular correction distributionally (indeed in the permitted
local $L^p$ topology).  On bounded-energy subsequences, the truncated
fields have weak $L^2(w)$ limits, identified by that correction.
It is this weighted convergence that passes the smeared divergence
equation, since test gradients belong to $L^2(w)$.
Adding back the singular correction recovers the unsmeared equation;
weighted Cauchy--Schwarz and the locally integrable singular flux give
$wE\in L^1_{\rm loc}$.  Distributional curl remains zero and the local
$L^p$ bounds persist, so the limit is a gradient in the prescribed
class.  These statements hold on compact sets of arbitrarily large
probability, by a summable allocation of the Markov bounds.
They therefore apply to the limiting law, not merely to a deterministic
subsequence of configurations.

Weighted quadratic energy is lower semicontinuous in this topology.
On a bounded cylinder this follows also from
\[
 \int w|F|^2
 =\sup_{\psi\in C_c^\infty}
       \left(2\langle F,\psi\rangle-\int w^{-1}|\psi|^2\right).
\]
Take the increasing supremum over cylinder heights for the full
extension integral.  Consequently a subsequential law $P$ satisfies
\eqref{eq:normalized-orbit-bound} with $P$ in place of $P_j$.
At this point the limit law is supported on compatible gradients and has
the required energy bounds; it may still have multiple charges.
It is horizontally stationary: for a fixed translation $z$, the
total variation discrepancy of an orbit law and its translate is at
most $|K_{S_j}\mathbin\triangle(K_{S_j}+z)|/S_j^2\to0$.

The $L^1$ pointwise and mean ergodic theorems along tempered
F\o lner sequences \citep[Section~1]{Lindenstrauss1999}, applied to
the centered square sequence in $\mathbb Z^2$, apply to each
integrable observable $\Gamma_{\eta_\ell}$.
The tiled sums of this observable are exactly the raw energies on integer squares;
positivity and nested squares extend the limits to every real radius.
Each limit is the conditional expectation of its observable on the
$\mathbb Z^2$-invariant sigma-algebra. It may depend on the realization;
no ergodicity of $P$ is assumed. Its expectation equals the mean of the
observable. Thus, on one event of full probability for all selected cutoffs,
\[
 a_{\eta_\ell}:=\mathbb E_P\mathcal W_{\eta_\ell}
       =\mathbb E_P\Gamma_{\eta_\ell}-\kappa g(\eta_\ell)
       \le m+2\kappa F_{\eta_\ell}.
\]
Fatou and the uniform lower bound give $\mathbb E_P\mathcal W\le m$.
This does not yet use minimality, because the limit could have
multiplicities.  It shows that $\mathcal W$ is finite and integrable.
Equation \eqref{eq:normalized-cutoff-bound} then gives dominated
convergence $a_{\eta_\ell}\to\mathbb E_P\mathcal W$.

\emph{Unit intensity and simplicity.}
The local count bound gives finite point intensity $\rho$.
In the logarithmic case the expected smeared PDE immediately gives
$\rho=1$ by horizontal stationarity.  In the extension case, let
$b_\eta(y)=\mathbb E_P(E_\eta)_y$ and let $\mu_{\eta,y}$ be the
vertical probability marginal of one smeared charge.  Jensen gives
$\int |y|^\gamma|b_\eta|^2<\infty$, and the mean PDE is
\[
 -\partial_y(|y|^\gamma b_\eta)
       =\kappa(\rho\mu_{\eta,y}-\delta_0).
\]
Outside $[-\eta,\eta]$ the weighted flux is constant on each tail.
They must vanish, since $\int_1^\infty y^{-\gamma}\,dy=\infty$.
Integrating the equation proves $\rho=1$.

For cutoffs $0<\alpha<\eta$ from the selected dyadic sequence, put
$d=\sum_p(f_\eta-f_\alpha)(\cdot-p)$, counting labeled particles.  Then
\begin{equation}
 w(|E_\alpha|^2-|E_\eta|^2)
  =\operatorname{div}[wd(E_\alpha+E_\eta)]
     +\kappa d(\nu_\alpha+\nu_\eta-2\delta_{\mathbb R^2}).
 \label{eq:normalized-two-cutoff-local}
\end{equation}
At positive cutoffs the kernels in $d$ are bounded and have finite
range.  The second count moment and Cauchy--Schwarz justify all local
expected products.  Test with a smooth horizontal function of integral
one and, if needed, a vertical cutoff equal to one on $|y|\le\eta$.
Stationarity annihilates the expected horizontal divergence; the
vertical cutoff varies outside the support of $d$.  Separating self,
distinct-labeled-pair, and background terms gives exactly
\begin{equation}
 D_{\alpha,\eta}=a_\alpha-a_\eta+2\kappa(F_\eta-F_\alpha),
 \qquad D_{\alpha,\eta}\ge0.
 \label{eq:normalized-two-cutoff-mean}
\end{equation}
Here $D_{\alpha,\eta}$ is the ordered-pair intensity of the kernel
\[
 \kappa\int(f_\eta-f_\alpha)(X-(p,0))
                (d\delta_q^{(\alpha)}+d\delta_q^{(\eta)})(X),
 \qquad p\ne q\text{ as labeled particles}.
\]
A coincident pair contributes $\kappa[g(\alpha)-g(\eta)]$.
If $c_{\rm coll}=\mathbb E_P\sum_{p\in Q}N_p(N_p-1)$, then
\[
 0\le\kappa[g(\alpha)-g(\eta)]c_{\rm coll}
       \le D_{\alpha,\eta}.
\]
The right side stays bounded as $\alpha\downarrow0$, whereas
$g(\alpha)\to\infty$.  Hence $c_{\rm coll}=0$; stationarity and a
countable tiling show simplicity almost surely.  Only now do we use
$\mathcal W\ge m$, proving $\mathcal W=m$ almost surely and
$a_{\eta_\ell}\to m$.

We have obtained simple minimizing fields. The remaining task is to select
one whose close pairs and vertical tail remain controlled during screening.

\emph{Vanishing defect and selection of a realization.}
Fatou in \eqref{eq:normalized-two-cutoff-mean}, with only
$\alpha\downarrow0$, gives
\begin{equation}
 0\le D_\eta:=\kappa\mathbb E_P
           \sum_{p\in Q}\sum_{q\ne p}q_\eta(p-q)
       \le m-a_\eta+2\kappa F_\eta\longrightarrow0.
 \label{eq:normalized-defect-small}
\end{equation}
The measure $\mu_\eta$ thus has finite mean intensity $D_\eta$.
The ergodic theorem therefore gives its spatial density
\[
 d_\eta(E)=\lim_{R\to\infty}\frac{\mu_\eta(K_R)}{R^2},
 \qquad \mathbb E_Pd_\eta=D_\eta
\]
almost surely.  For $k=1$, apply the same ergodic theorem
to the raw energy over each fixed integer height $|y|>z>1$.
These integrable observables decrease to zero, so their spatial
densities decrease to zero almost surely.  Intersect all the energy,
tail, and selected defect events, a countable family of full-measure
events. For each selected $\eta$, choose a realization $(E,\mathcal C)$
in this intersection with $d_\eta\le2D_\eta$; if $D_\eta=0$, choose
$d_\eta=0$.  This realization may depend on $\eta$.
For every fixed $\varepsilon>0$, define
\begin{equation}
 e_{\varepsilon,R}:=\frac1{\varepsilon^4R^2}
 \int_{K_R\times\{|y|>\frac14\varepsilon^2R\}}w|E|^2
                         \longrightarrow0.
 \label{eq:normalized-moving-tail}
\end{equation}
The coefficient $\frac14$ in this threshold is fixed independently of
the cutoffs and the realization.
Indeed, compare the moving height with each fixed integer height and
then let that fixed height tend to infinity.  No uniform rate in
$\varepsilon$ is needed.

\end{proof}

\subsection{Screening at a fixed cutoff}

We use the rectangular shell construction of
Petrache and Serfaty~\citep[Section~6.2]{PetracheSerfaty2017}. The next lemma records the
precise input and output needed here; the proof supplies the local field
equations and their constants in the present normalization.

\begin{lemma}[Screening without new close pairs]\label{lem:field-screening}
Fix the auxiliary cutoff $\eta'$ chosen above and $0<\eta<\eta'$.
Let $E$ be a compatible simple gradient field for which
\[
 M_{R,\rho}:=R^{-2}Q_\rho(E,R)
\]
has a finite limit for $\rho=\eta,\eta'$. If $k=1$, suppose also that
$e_{\varepsilon,R}$ defined in \eqref{eq:normalized-moving-tail} tends to
zero for every fixed $\varepsilon>0$; set $e_{\varepsilon,R}=0$ if $k=0$.
For every $0<\varepsilon<1/8$ and every sufficiently large integer $R$,
there are exactly $R^2$ distinct points $\widehat\Lambda\subset K_R$
and a field $\widehat E$ with zero exterior normal flux satisfying
\[
 -\operatorname{div}(w\widehat E)
 =\kappa\left(\sum_{p\in\widehat\Lambda}\delta_{(p,0)}
                 -\mathbf1_{K_R}\delta_{\R^2}\right).
\]
Its truncation has finite weighted energy and obeys
\begin{align*}
 R^{-2}\int_{K_R\times\R^k}w|\widehat E_\eta|^2
 &\le M_{R,\eta}
 +C_{s,\eta'}\varepsilon(1+g(\eta))
                     (1+M_{R,\eta}+M_{R,\eta'})+o_R(1)\\
 &\quad+C_s e_{\varepsilon,R}
                     \varepsilon^{1+\min(2\gamma,0)}.
\end{align*}
Every pair in $\widehat\Lambda$ at distance below $2\eta$ is a retained
pair of the original configuration in $K_R$, and every point has distance
at least $a_0/2$ from $\partial K_R$. The field $\widehat E$ need not be
a gradient. All constants are independent of $R$ and the realization;
$R$ is chosen after $\eta',\eta$, the field, and $\varepsilon$.
\end{lemma}

\begin{proof}
Keep the field and both cutoffs fixed. We first choose the inner cylinder
where its truncated field is retained, then construct the shell and caps.
Coarea supplies an inner rectangle $K'_R$ whose corresponding faces
are at distances between $\varepsilon R$ and $2\varepsilon R$ from
the outer boundary, with
\[
 \int_{\partial K'_R\times\mathbb R^k}w|E_\eta|^2
       \le C M_{R,\eta}\varepsilon^{-1}R,
 \quad
 \int_{(\partial K'_R)_b\times\mathbb R^k}w|E_{\eta'}|^2
       \le C M_{R,\eta'}\varepsilon^{-1}R,
\]
where $b$ is a fixed sufficiently large shell width and
$M_{R,\rho}=R^{-2}Q_\rho(E,R)$.
For $k=1$, choose
$\ell\in[\frac12\varepsilon^2R,\varepsilon^2R]$.
Such a height with the required trace bound exists because
\begin{align*}
 \int_{\frac12\varepsilon^2R}^{\varepsilon^2R}
   \left(\int_{K'_R\times\{\pm t\}}w|E|^2\right)\,dt
 &\le \int_{K_R\times\{|y|>\frac14\varepsilon^2R\}}w|E|^2\\
 &=\varepsilon^4R^2e_{\varepsilon,R}.
\end{align*}
Dividing by the interval length $\frac12\varepsilon^2R$ gives
\[
 \int_{K'_R\times\{\pm\ell\}}w|E|^2
                    \le2\varepsilon^2R e_{\varepsilon,R}.
\]
For $k=0$, retain the choice $\ell=\varepsilon^2R$.
In both cases $R$ is then taken large enough that $\ell>2\sqrt2$.
The fixed-cutoff energy limits bound both $M$'s as $R\to\infty$.

Let $D_0=K'_R\times[-\ell,\ell]^k$, $A_{\rm sh}=|K_R|-|K'_R|$, and,
when $k=1$,
let $T_\pm$ be the outward weighted flux of the old field through the
top and bottom of $D_0$.  The cap constants must be
\begin{equation}
 C^\pm=-\frac{T_\pm}{A_{\rm sh}\ell^\gamma}.
 \label{eq:normalized-two-caps}
\end{equation}
They impose separate compatibility in the disconnected upper and
lower caps $K_R\times(\ell,R)$ and $K_R\times(-R,-\ell)$.
For $k=0$ all $T_\pm,C^\pm$ and cap terms are omitted.
Let $\Lambda_0=\{p\in\mathcal C:\operatorname{dist}(p,\partial K'_R)\le\eta\}$
be the retained crossing centers and
$\rho_0=\sum_{p\in\Lambda_0}N_p\delta_p^{(\eta)}$ their full
smeared measure.  The retained old centers are
$\Lambda_{\rm old}=(\mathcal C\cap K'_R)\cup\Lambda_0$;
$N_{\rm old}$ counts all of them, including multiplicity if present.
All pieces assembled below are finite-energy components of the
\emph{truncated} field.  In particular the new-charge component uses
smeared sources; the point singularities are restored only after the
pieces have been assembled.
On a shell cell $H_i$, denote the old side flux by
$B_i=\int_{(\partial K'_R\cap\partial H_i)\times[-\ell,\ell]^k}wE_\eta\cdot\nu_0$,
where $\nu_0$ points outward from $D_0$.
The truncated shell field is the sum of five components, constructed below:
\begin{center}
\begin{tabular}{@{}ll@{}}
\toprule
Field & Role in the shell\\
\midrule
$E_{\rm cross}$ & Complete the retained smeared charges across the inner boundary.\\
$E_{\rm vert}$ & Match the upper and lower cap data.\\
$E_{\rm side}$ & Match the old lateral flux and the ideal background density.\\
$E_{\rm round}$ & Correct the density when cell masses are rounded to integers.\\
$E_{\rm new}$ & Place new smeared charges against the rounded background.\\
\bottomrule
\end{tabular}
\end{center}
For $k=0$, $E_{\rm vert}=0$ and there are no caps.
We use two independent fixed partitions of the shell, so that corner
completion and integer charge counts do not require moving its faces.

For the large cells, include the coordinates of every inner and outer
face as breakpoints in each coordinate axis.  Subdivide each resulting
interval into lengths in $[\ell,2\ell]$, and remove the product cells
inside $K'_R$.  The remaining rectangles $H_i$ form a connected annular
grid.  Adjacent rectangles share full edges of length at least $\ell$,
and there are $O_\varepsilon(1)$ rectangles at fixed $\varepsilon$.
All the one-dimensional intervals exceed $\ell$ for sufficiently large
$R$ and the chosen small $\varepsilon$.

We spell out the field estimates needed in this construction.
First consider a fixed cell $\Omega=S\times[-1,1]^k$, where $S$ is a
rectangle with sidelengths in $[1,2]$; when $k=0$, take $\Omega=S$.
Suppose the cell carries a restriction $\nu$ of smeared charges, of total
mass $Q$. A constant outward normal component $b_F$ on a chosen lateral
face $F$ is compatible with $Lh=\kappa\nu$ precisely when
\[
                 b_F\int_Fw=-\kappa Q.
\]
Use the representative $u$ of weighted cell mean zero; compatibility
allows this choice because the Neumann functional annihilates constants.
The weighted Poincar\'e and trace estimates on these cells can be seen
directly. Apply ordinary Poincar\'e and trace in the horizontal variables,
then integrate against $|y|^\gamma$. In the vertical variable use
\[
 |v(y)-v(z)|^2
 \le\left(\int_{-1}^1|t|^\gamma|v'(t)|^2\,dt\right)
     \left(\int_{-1}^1|t|^{-\gamma}\,dt\right).
\]
Both $w$ and $w^{-1}$ are integrable on this interval because
$-1<\gamma<1$. Averaging this inequality proves weighted Poincar\'e
and the trace bound at $y=0$; horizontal trace gives the lateral-face
bound. The constants are uniform over the stated sidelengths.

For the extension, reflect across the coordinate faces of the base rectangle $S$, which
preserves $w$. If $k=1$, extend in $y$ only to $[-3/2,3/2]$, reflecting
each added collar into $1/2<|y|<1$, where the original and reflected
weights are comparable. Multiply by a smooth cutoff equal to one on
$\Omega$ and supported in this fixed enlargement. Poincar\'e controls
the term $u\nabla\chi$ from that multiplication, giving an extension
$Tu$ with
\[
 \|\nabla Tu\|_{L^2(w)}
       \le C_\gamma\|\nabla u\|_{L^2(w;\Omega)}.
\]
When $k=0$, the same construction uses only the horizontal reflections.

Put $g_\eta=g-f_\eta$, the fundamental solution capped at $g(\eta)$.
The retained centers contributing to this cell lie in a fixed enlargement
of it. On the support of $Tu$, radial integration of the fundamental
solution's flux gives
\[
 \int_{\operatorname{supp}(Tu)}w
                  |\nabla g_\eta(X-(p,0))|^2\,dX
       \le C_s(1+g(\eta)).
\]
For $s>0$ the full-space integral is $\kappa g(\eta)$; for the
logarithm the integral over a ball of fixed radius is
$\kappa\log(1/\eta)+O(1)$. Only this local integral is used.
Positivity of the smearing measure then gives
\[
 \left|\int_\Omega u\,d\delta_p^{(\eta)}\right|
 \le\int |Tu|\,d\delta_p^{(\eta)}
 =\kappa^{-1}\left|\int w\nabla g_\eta(\cdot-p)
                                      \cdot\nabla|Tu|\right|
 \le C\sqrt{1+g(\eta)}\,\|\nabla u\|_{L^2(w)}.
\]
The Neumann functional annihilates constants.  Lax--Milgram therefore
gives energy at most $C N^2(1+g(\eta))$ when at most $N$ labeled
charges contribute and the sum of the absolute prescribed face fluxes
is at most $C\kappa N$.  The same proof allows several constant face
data whose total is $-\kappa Q$: apply the trace bound to each face.
This proves the local estimate needed from
Petrache and Serfaty~\citep[Lemma~6.6]{PetracheSerfaty2017}, with the outward flux fixed by
$Lh=\kappa\nu$.

For crossing-charge completion, form a separate tensor grid with
sidelengths in $[1,2]$, again containing all inner and outer face
coordinates.  Restrict $\rho_0$ to its exterior cells times
$[-1,1]^k$.  Its exterior support lies within distance $2\eta$ of
$\partial K'_R$.  With $2\eta<1/2$, every such cell either lies in
the first layer outside exactly one inner coordinate face, and shares
a full face with $K'_R$, or is one of the four diagonal corner cells.
Route each corner cell to an adjacent first-layer cell through their
full shared face.  These are trees of uniformly bounded size; cells
not receiving a corner contribution are single-cell trees.

All face data in the following construction are constant normal
derivatives: a prescribed integrated weighted flux $F_0$ means the value
$F_0/\int_Fw$ on that face. On a leaf, prescribe outward weighted flux $-\kappa$ times its
restricted source mass through the face leading to the root.  At a
root, prescribe the opposite incoming flux through each child face,
and export $-\kappa$ times the total tree mass through its full inner
face.  All other face data are zero.  More generally the outgoing
flux of a tree vertex is $-\kappa$ times the mass in its subtree.
The sum of its outward fluxes is exactly $-\kappa$ times its own
source mass, so the cell problems are compatible.  Internal normal
fluxes cancel.  This gives one global completion field $E_{\rm cross}$
with source $\kappa\rho_0$ on the exterior part, and a nonnegative
outlet density
\[
 j=-E_{\rm cross}\cdot\nu_{\rm shell}
   =\sum_{F\ {\rm outlet}}\mathbf1_F
                   \frac{\kappa Q_F}{\int_Fw}
       \quad\hbox{on the inner shell boundary}.
\]
Here $Q_F$ is the total tree mass and $\nu_{\rm shell}=-\nu_0$.
Figure~\ref{fig:normalized-corner-routing} shows the local corner relay.

\begin{figure}[tb]
\centering
\begin{tikzpicture}[
  x=1.8cm,y=1.8cm,
  >=Latex,
  line cap=round,
  font=\small
]
  \fill[black!7] (0,0) rectangle (1.5,1.5);
  \draw[black!45,line width=.4pt] (-1.5,-1.5) rectangle (1.5,1.5);
  \draw[black!45,line width=.4pt] (0,-1.5) -- (0,0);
  \draw[line width=.65pt] (-1.5,0) -- (0,0);
  \draw[line width=1pt] (0,1.5) -- (0,0) -- (1.5,0);

  \node at (.83,1.05) {$K'_R$};
  \node at (-.85,1.26) {side root};
  \node at (-.78,-1.12) {corner cell};
  \node[align=center] at (.80,-.98) {another\\side cell};

  \draw[->,line width=.85pt] (-.82,-.65) -- (-.82,.52)
    node[pos=.42,left=3pt] {$Q_{\rm d}$};
  \draw[->,line width=.85pt] (-.74,.68) -- (0,.68);
  \node at (-.75,.95) {$Q_{\rm d}+Q_{\rm r}$};

  \draw[line width=.65pt] (.08,.08) circle[radius=.22];
  \fill (.08,.08) circle[radius=1.25pt];
  \draw[black!65,line width=.4pt] (.11,.11) -- (.43,.34);
  \node[anchor=west,inner sep=2pt] at (.43,.34) {$p$};
  \draw[black!65,line width=.4pt] (.25,-.06) -- (.53,-.29);
  \node[anchor=west,inner sep=2pt] at (.53,-.29) {$\eta$-ball};
\end{tikzpicture}
\caption{Auxiliary-cell routing near an inner corner (base-plane view).
The corner cell has no full face on $K'_R$, so its restricted
smeared mass $Q_{\rm d}$ is routed to the side root, whose own restricted
mass is $Q_{\rm r}$. The arrows indicate mass routing. On each outgoing
face, the assigned integrated outward flux is $-\kappa$ times the indicated
mass; the internal face fluxes cancel.
One retained center and the section of its $\eta$-ball are illustrated.
These auxiliary cells are independent of the large cells $H_i$ and of
the rectangles used for new particles. The completion field
$E_{\rm cross}$ is added to the other shell fields.}
\label{fig:normalized-corner-routing}
\end{figure}

The completion field is not cut at the boundaries of the large cells
$H_i$; thus no trace on an arbitrary such cut is required.
The local count estimate and bounded overlap of the unit trees give
\[
 \int w|E_{\rm cross}|^2
       \le C(1+g(\eta))\sum_{\rm trees}N_{\rm tree}^2,
 \qquad
 \int_{\partial K'_R\times\mathbb R^k}w j^2
       \le C\sum_{\rm trees}N_{\rm tree}^2
       \le C(1+M_{R,\eta'})\varepsilon^{-1}R.
\]
The additive $1$ in the count-square bound is required.

The crossing field now completes every retained smeared charge and provides
its actual outlet flux on the inner boundary. We next balance that flux
against the background, then choose an integer number of new particles.

For each large cell set
$Q_i=\int_{(\partial K'_R\cap\partial H_i)\times[-\ell,\ell]^k}w j$,
its actual inner-face outlet flux.  This need not equal the restricted
source mass $\kappa\rho_0(H_i\times[-\ell,\ell]^k)$, because the
completion field may transport that source across a large-cell boundary.
Define the ideal density by
\begin{equation}
 \kappa(m_i-1)|H_i|
       =B_i-\ell^\gamma(C^++C^-)|H_i|-Q_i.
 \label{eq:normalized-shell-density}
\end{equation}
Since $\sum_iQ_i=\kappa\rho_0(D_0^c)$, Gauss' equation and
\eqref{eq:normalized-two-caps} give
$\sum_i m_i|H_i|=R^2-N_{\rm old}=:N_{\rm new}\in\mathbb N$.
At fixed parameters the lateral-flux, cap-flux and outlet bounds make
$\max_i|m_i-1|\to0$ as $R\to\infty$; choose this maximum below $1/4$.
For example the outlet term divided by $|H_i|$ is bounded by
$C(1+M_{R,\eta'})\varepsilon^{-5}R^{-1}$.

Set $r_i=m_i|H_i|$.  Round the $r_i$ down, and increase the required
number of entries by one, obtaining integers $n_i$ with
$\sum_i n_i=N_{\rm new}$ and $|n_i-r_i|\le1$.
Put $\widetilde m_i=n_i/|H_i|$.  For large $R$ these densities lie
in $[1/2,3/2]$.  To compensate this rounding, write
$\delta_i=n_i-r_i$, whose sum is zero, and choose a rooted spanning
tree of the large-cell adjacency graph.  On the face from a child
cell $i$ to its parent prescribe outward weighted flux
$-\kappa\sum_{a\text{ in the subtree of }i}\delta_a$, with opposite
flux on the adjacent cell.  The total outward flux of cell $i$ is
$-\kappa\delta_i$.  There is therefore a compatible Neumann field
$E_{{\rm round},i}$ with source
$\kappa\delta_i|H_i|^{-1}\delta_{\mathbb R^2}$, these face fluxes,
and zero flux on all other faces, including the caps.
The paired fluxes make the union $E_{\rm round}$ a global field.
It has no flux through the inner or outer shell boundary.
After dilation by $\ell$, weighted trace and Poincar\'e on rectangles
of bounded aspect ratio give
\[
 \int w|E_{\rm round}|^2\le
 \begin{cases}C_{s,\varepsilon}\ell^{-s},&s>0,\\
 C_\varepsilon,&s=0.
 \end{cases}
\]
Indeed there are $O_\varepsilon(1)$ cells and each tree flux is
$O_\varepsilon(1)$; the unit-scale variational problem has bounded
energy, and dilation scales it by $\ell^{-s}$ for Riesz kernels and
by one for the logarithm.  This is $o(R^2)$ at fixed parameters.
For sufficiently large $R$, every sidelength of $H_i$ is at least
$\ell>2\sqrt2\ge2\widetilde m_i^{-1/2}$, and
$\widetilde m_i|H_i|=n_i\in\mathbb N$.
Thus Petrache and Serfaty~\citep[Lemma~6.3]{PetracheSerfaty2017} partitions $H_i$ into
$n_i$ rectangles of area $1/\widetilde m_i$, with every sidelength
in
\[
 \left[\frac14\widetilde m_i^{-1/2},
                      4\widetilde m_i^{-1/2}\right]
                  \subset[a_0,a_1].
\]
In particular, the lower bound $a_0$ fixed at the start applies to
every new-charge rectangle, independently of all later choices.

The vertical correction on each shell cell is the gradient of
\[
 h_2(y)=\frac{\ell^\gamma}{1-\gamma}
 \begin{cases}C^+y^{1-\gamma},&y\ge0,\\
 C^-|y|^{1-\gamma},&y\le0.
 \end{cases}
\]
It has source $-\ell^\gamma(C^++C^-)\delta_{\mathbb R^2}$.
The side-flux correction has the opposite cap source plus
$\kappa(m_i-1)\delta_{\mathbb R^2}$ and datum
$-E_\eta\cdot\nu_0+j$ on the
inner side, with zero data elsewhere.  Equation
\eqref{eq:normalized-shell-density} is exactly its compatibility.

The large-cell estimate used for this correction is as follows.
Let $H$ have sidelengths in $[\ell,2\ell]$, set
$\Omega=H\times[-\ell,\ell]$ and
$\Sigma=\partial H\times[-\ell,\ell]$, and let $-1<\gamma<1$.
For a constant $a$ and $b\in L^2(\Sigma,|y|^\gamma)$ satisfying
\[
              a|H|=-\int_\Sigma |y|^\gamma b,
\]
the weak Neumann problem
\[
 -\operatorname{div}(|y|^\gamma\nabla h)
                  =a\delta_{\mathbb R^2}\quad\hbox{in }\Omega,
 \qquad \partial_\nu h=b\quad\hbox{on }\Sigma,
 \qquad \partial_\nu h=0\quad\hbox{on }H\times\{\pm\ell\}
\]
has a solution modulo constants and obeys
\begin{equation}
 \int_\Omega |y|^\gamma|\nabla h|^2
 \le C_\gamma\left(a^2\ell^{3-\gamma}
              +\ell\int_\Sigma |y|^\gamma|b|^2\right)
 \le C_\gamma\ell\int_\Sigma |y|^\gamma|b|^2.
 \label{eq:large-cell-side-estimate}
\end{equation}
The constants depend only on $\gamma$ and the stated aspect-ratio
bounds.  To see the scaling, choose the weighted-mean-zero
representative.  Weighted Poincar\'e and the one-dimensional trace
inequalities, after dilation by $\ell$, give
\[
 \int_H|h(x,0)|^2\le C_\gamma\ell^{1-\gamma}
                              \int_\Omega |y|^\gamma|\nabla h|^2,
 \qquad
 \int_\Sigma |y|^\gamma|h|^2
          \le C_\gamma\ell\int_\Omega |y|^\gamma|\nabla h|^2.
\]
Testing the weak equation gives the first inequality in
\eqref{eq:large-cell-side-estimate}.
Compatibility, $|H|\asymp\ell^2$, and
$\int_\Sigma|y|^\gamma\asymp\ell^{\gamma+2}$ give
$a^2\ell^{3-\gamma}\le C_\gamma\ell\int_\Sigma|y|^\gamma|b|^2$,
which proves the second.  This is the constant-source case of
Petrache and Serfaty~\citep[Lemma~6.4]{PetracheSerfaty2017}, with the source coefficient
written explicitly.  For the logarithm use $\Omega=H$,
$-\Delta h=a$, and $\partial_\nu h=b$ on $\partial H$;
ordinary Poincar\'e, trace, and compatibility give the same final
bound $C\ell\int_{\partial H}|b|^2$.

For the side field on $H_i$, take
\[
 a_i=\kappa(m_i-1)+\ell^\gamma(C^++C^-),\qquad
 b_i=-E_\eta\cdot\nu_0+j
       \quad\hbox{on its inner face,}
\]
with $b_i=0$ on its other lateral faces.  The cap term is omitted
when $k=0$.  The density equation gives
$a_i|H_i|=B_i-Q_i=-\int_{\Sigma_i}w b_i$ exactly, so the cap
contribution to the constant source is included in the compatibility
estimate.  The inner faces partition the old lateral boundary.
Consequently the old-field coarea bound and the outlet bound yield
\begin{align*}
 \int w|E_{\rm side}|^2
 &\le C_s\ell\sum_i\int_{\Sigma_i}w|b_i|^2\\
 &\le C_s\ell\left(
      \int_{\partial K'_R\times[-\ell,\ell]^k}w|E_\eta|^2
          +\int_{\partial K'_R\times\mathbb R^k}wj^2\right)\\
 &\le C_{s,\eta'}(\varepsilon^2R)(\varepsilon^{-1}R)
                    (M_{R,\eta}+1+M_{R,\eta'})\\
 &=C_{s,\eta'}\varepsilon R^2(M_{R,\eta}+1+M_{R,\eta'}).
\end{align*}
Here $\Sigma_i=\partial H_i\times[-\ell,\ell]^k$.
This calculation supplies the length factor in the side-field row
of the shell-energy estimate.  Its constant may depend on the single
fixed auxiliary cutoff $\eta'$, but not on $\eta$, $\varepsilon$,
$R$, or the realization.

Finally, let $S\subset H_i$ be one of the area-$1/\widetilde m_i$
rectangles, with barycenter $p$, and work on the fixed cylinder
$\Omega_S=S\times(-1,1)^k$ (or $\Omega_S=S$ when $k=0$).
Take the $\eta$-truncation of the neutral Neumann field with source
$\kappa(\delta_{(p,0)}-\widetilde m_i\mathbf1_S\delta_{\mathbb R^2})$
and zero normal component on its entire boundary.  The truncated source is
$\kappa(\delta_p^{(\eta)}-\widetilde m_i\mathbf1_S\delta_{\mathbb R^2})$.
Extend this truncated field by zero into $H_i\times[-\ell,\ell]^k$
and the shell.  The sidelength bounds $[a_0,a_1]$, fixed height, and
$\eta<a_0/16$ keep the charge uniformly away from the cell boundary.
To verify its energy, write the point-source potential as
$g(\cdot-(p,0))+v$. The correction solves
\[
 Lv=-\kappa\widetilde m_i\mathbf1_S\delta_{\R^2},\qquad
 \partial_\nu v=-\partial_\nu g(\cdot-(p,0))
                       \quad\hbox{on }\partial\Omega_S.
\]
The weighted trace and Poincar\'e estimates give a uniform bound on
$\int_{\Omega_S}w|\nabla v|^2$; no smoothness across $y=0$ is needed.
Test this compatible weak equation against $g_\eta(\cdot-(p,0))$.
The resulting plane integral is uniformly bounded because $s<2$, and
the boundary terms are uniformly bounded by the clearance from the pole.
Hence $\int_{\Omega_S}w\nabla v\cdot\nabla g_\eta=O_s(1)$.
The capped fundamental solution itself has energy
$\kappa g(\eta)+O_s(1)$ in $\Omega_S$, so the truncated cell field has
the same energy, uniformly over the allowed rectangles.
In addition to the preceding $o(R^2)$ bound for $E_{\rm round}$,
the shell fields satisfy
\begin{align*}
 \int w|E_{\rm cross}|^2
   &\le C(1+g(\eta))(1+M_{R,\eta'})\varepsilon^{-1}R,\\
 \int w|E_{\rm vert}|^2
   &\le C e_{\varepsilon,R}\varepsilon^3R^2,\\
 \int w|E_{\rm side}|^2
   &\le C_{s,\eta'}\varepsilon(M_{R,\eta}+1+M_{R,\eta'})R^2,\\
 \int w|E_{\rm new}|^2
   &\le C(1+g(\eta))N_{\rm new},
\end{align*}
where $N_{\rm new}=O(\varepsilon R^2)+o(R^2)$ at fixed parameters.
The first bound uses the sum of squared tree counts; the third
uses the weighted rectangle trace estimate with height $\ell$.
Write $\nu_{{\rm new},\eta}=\sum_{p\ {\rm new}}\delta_p^{(\eta)}$.
The five shell fields have combined source
\[
 \kappa\rho_0
 -\ell^\gamma(C^++C^-)\delta_{\mathbb R^2}
 +\bigl[\kappa(m_i-1)+\ell^\gamma(C^++C^-)\bigr]\delta_{\mathbb R^2}
 +\kappa(\widetilde m_i-m_i)\delta_{\mathbb R^2}
 +\kappa(\nu_{{\rm new},\eta}-\widetilde m_i\delta_{\mathbb R^2}),
\]
which is $\kappa(\rho_0+\nu_{{\rm new},\eta}-\delta_{\mathbb R^2})$.
On the inner shell boundary the crossing and side data sum to
$-j+(-E_\eta\cdot\nu_0+j)=-E_\eta\cdot\nu_0$.
The rounding field has zero inner data and paired internal data;
the completion field remains global across the large-cell boundaries.
Thus all interface fluxes cancel.  We sum squared energies only in
the shell, where a fixed factor of five is harmless.  The old field
inside $D_0$ is retained exactly, with coefficient one in the energy
bound; no global factor multiplying its energy is introduced.

The shell now has the required source and lateral flux. It remains to close
the two vertical ends without changing these sources.

The upper and lower weighted harmonic caps are solved separately.
On the inner horizontal face of the upper or lower cap, prescribe the
outward normal derivative $-E_\eta(x,\pm\ell)\cdot(\pm e_y)$ for
$x\in K'_R$, and $-C^\pm$ for $x\in K_R\setminus K'_R$; here $e_y$
is the positive vertical unit vector.  Prescribe zero normal derivative
on every other face.  By \eqref{eq:normalized-two-caps}, the weighted
integral of these data is
$-T_\pm-A_{\rm sh}\ell^\gamma C^\pm=0$, so each cap problem is
compatible.  If $u$ is adjusted by its ordinary mean on one cap,
ordinary trace and Poincar\'e, followed by weight comparison, give
\[
 \int_{K_R\times\{\ell\}}w|u|^2
   \le CR\max\{1,(\ell/R)^\gamma\}
                              \int_{\rm cap}w|\nabla u|^2.
\]
Testing the compatible cap problem yields total cap energy at most
\[
       C e_{\varepsilon,R}
                     \varepsilon^{1+\min(2\gamma,0)}R^2.
\]
The extra factor for $\gamma<0$ is retained.  It vanishes at fixed
$\varepsilon$ by \eqref{eq:normalized-moving-tail}.
There are no cap terms for $k=0$.

Retain $E_\eta$ in $D_0$, use the sum of the five finite-energy
components in the shell, and use the two harmonic cap fields above
and below.  Extend this assembled field $\widehat F_\eta$ by zero past the
outer square and the cap heights.  It has zero exterior normal flux
and source $\kappa(\sum_{p\in\widehat\Lambda}\delta_p^{(\eta)}
-\mathbf1_{K_R}\delta_{\mathbb R^2})$, where
$\widehat\Lambda$ consists of the retained old and new barycenter
charges, exactly $R^2$ points.  Now define
\[
 \widehat E=\widehat F_\eta+
               \sum_{p\in\widehat\Lambda}\nabla f_\eta(\cdot-(p,0)).
\]
The singular corrections lie inside the outer square by its clearance.
Thus $\widehat E_\eta=\widehat F_\eta$ exactly, and the fundamental-solution
identity proves the required unsmeared divergence equation, with
zero exterior normal flux.  At fixed cutoffs the preceding bounds
give
\begin{align}
 R^{-2}\int_{K_R\times\mathbb R^k}w|\widehat E_\eta|^2
 &\le M_{R,\eta}
   +C_{s,\eta'}\varepsilon(1+g(\eta))
                 (1+M_{R,\eta}+M_{R,\eta'})+o_R(1)\notag\\
 &\quad+C_s e_{\varepsilon,R}
                       \varepsilon^{1+\min(2\gamma,0)}.
 \label{eq:normalized-screen-estimate}
\end{align}
Here $R$ is chosen after $\eta,\eta',\varepsilon$.
For example, the count estimate is controlled once $R\gg(1+M_{R,\eta'})\varepsilon^{-5}$.

Every new barycenter is at least $a_0/2$ from each face of its
rectangle.  Retained old points lie inside $K'_R$ or within $\eta$
of its boundary.  Since $\eta<\eta_0\le a_0/16$ was fixed before
the realization was selected, every new--old distance is at least
$a_0/2-\eta>2\eta$.
Choose the integer radius large enough that
$\varepsilon R-\eta\ge a_0/2$; this gives the same clearance at
the outer boundary and hence at every reflected seam.
New--new distances and all reflection-seam distances exceed $2\eta$.
Thus every close pair after screening was already an old close pair.

\end{proof}

\subsection{Reflection, projection, and removal of the cutoff}

\begin{proof}[Proof of Proposition~\ref{prop:normalized-cubic-approximation}]
Lemma~\ref{lem:field-cutoff-control} proves existence and the uniform lower
bound of the cutoff limit, as well as finiteness of $m$. To approximate $m$,
choose a selected dyadic cutoff $\eta$ and a realization supplied by
Lemma~\ref{lem:stationary-field-minimizers}. Apply
Lemma~\ref{lem:field-screening} after choosing $\varepsilon$, and then
choose its integer radius $R$. We now convert that screened field into a
canonical periodic gradient and estimate its energy after removing the
cutoff.

Reflect the screened square in both coordinate directions and repeat
with period $2R\mathbb Z^2$.  The reflected vector components have
matching zero normal flux, so there are no interface charges.  The
period cell has $4R^2$ points and area $4R^2$.
Reflection multiplies both energy and defect by four.  The nonnegative
defect of the resulting configuration is at most the old anchored
defect in $K_R$, divided by $R^2$, and hence tends to at most
$d_\eta(E)\le2D_\eta$.

Write $E^{\rm per}$ for the reflected field, which need not be curl free,
and let $E^{\rm can}$ be the canonical periodic gradient with exactly
its charges and background. At fixed cutoff,
$D=E^{\rm per}-E^{\rm can}=E^{\rm per}_\eta-E^{\rm can}_\eta$
belongs to $L^2(w)$ on the full period cylinder and has zero weighted
divergence.  Test against the canonical truncated potential, with
vertical cutoffs tending to infinity; its nonconstant Fourier modes
decay vertically and its zero mode has compactly supported derivative.
The cross term vanishes, so
\[
 \int_{T\times\mathbb R^k}w|E^{\rm per}_\eta|^2
  =\int_{T\times\mathbb R^k}w|E^{\rm can}_\eta|^2
                       +\int_{T\times\mathbb R^k}w|D|^2.
\]
Canonical projection therefore decreases the cutoff energy and leaves
the configuration defect unchanged. It also returns to the prescribed
admissible class. For $s>0$, finite truncated $L^2(w)$ energy gives local
unweighted $L^p$ by weighted H\"older when $p<\min(2,2/s)$, and each of
the finitely many local singular corrections belongs to $L^p$ when
$p<3/(s+1)$.
For the logarithm both bounds reduce to $p<2$. Weighted Cauchy--Schwarz
and the locally integrable singular flux give $wE^{\rm can}\in
L^1_{\rm loc}$. Thus the canonical gradient is a member of
$\mathcal A_1$ with the exponent fixed at the beginning, even though the
intermediate screened field was not required to be a gradient.

For a finite simple periodic configuration, integrate
\eqref{eq:normalized-two-cutoff-local} over its period cell and let the
auxiliary cutoff tend to zero. This gives
\begin{equation}
 V(\mathcal C)-\mathcal W_\eta(E^{\rm can}_{\mathcal C})
   =\frac{\kappa}{|T|}\sum_{p\in T}\sum_{q\ne p}q_\eta(p-q)
                        -2\kappa F_\eta.
 \label{eq:normalized-periodic-defect}
\end{equation}
The second sum includes periodic images; the self image at zero is
the only omitted particle.

Keep the auxiliary cutoff $\eta'$ fixed at its initial value as $\eta$
tends to zero.  For each selected $\eta$, take the realization chosen
above, and choose $\varepsilon$ so the middle
term of \eqref{eq:normalized-screen-estimate} is as small as desired;
then choose an arbitrarily large integer $R$ so that all energy and
defect limits, the moving-tail estimate, the size conditions, and the
seam clearance hold simultaneously.  Subtract $\kappa g(\eta)$,
project canonically, and apply \eqref{eq:normalized-periodic-defect}.
Equations \eqref{eq:normalized-cutoff-bound} and
\eqref{eq:normalized-defect-small} give
\[
 V(\mathcal C_{\eta,R})
   \le m+2\kappa F_\eta+o(1)+2D_\eta-2\kappa F_\eta
   =m+2D_\eta+o(1).
\]
Let the chosen cutoffs tend to zero, taking the screening thickness
and integer radius in that order for each cutoff.  The canonical
periodic gradients belong to $\mathcal A_1$, so their values are also
at least $m$.  This proves the asserted cubic minimizing sequence.
Every use of the box limit preceded removal of the cutoff; no
interchange of the two limits has been made.
\end{proof}
